\ProvidesFile{For.tex}[Compile-time цикл for]


\subsection{Compile-time цикл for}

\subsubsection{Что мы хотим}

Давайте в начале обозначим проблему, которую мы хотим решать.
Предположим, у нас есть следующая шаблонная функция:
\begin{cppcode}
template<int index>
void function(double x, char y) {...}
\end{cppcode}
И мы хотим вызвать эту функцию несколько раз так:
\begin{cppcode}
void function<0>(1.2, 'c');
void function<1>(1.2, 'c');
void function<2>(1.2, 'c');
void function<3>(1.2, 'c');
void function<4>(1.2, 'c');
\end{cppcode}
Чтобы не заниматься копипастой, хотелось бы иметь возможность подобный вызов сделать так:
\begin{cppcode}
For<From{0}, To{5}>::Do<function>(1.2, 'c');
\end{cppcode}
К сожалению в точности так не получится, просто потому что нельзя передать шаблон функции в качестве параметра шаблона.
Но можно функцию завернуть внутрь шаблонной структуры, которую уже можно легко передать в качестве параметра.

\subsubsection{Имплементация \texttt{For}}

В начале заведем несколько служебных псевдонимов:
\begin{cppcode}
using From = AliasE<int, struct FromForLoopTag>;
using To = AliasE<int, struct ToForLoopTag>;
using Step = AliasE<int, struct StepForLoopTag>;
\end{cppcode}
Здесь используются strong type aliases из раздела~\ref{section::CtorArgs}.
Тип \verb"int" используется потому что целочисленные литералы имеют такой тип.
Знаковый тип для лучшей поддержки арифметики над значениями, чтобы не мучиться с арифметикой и переполнением в отрицательную сторону.
А псевдонимы используются для того, чтобы был более ясный интерфейс, сравните
\begin{cppcode}
For<0, 5>::Do<function>(1.2, 'c');
For<From{0}, To{5}>::Do<function>(1.2, 'c');
\end{cppcode}
Может быть стоило использовать \verb"Begin" и \verb"End", ну да ладно.

Теперь как имплементировать такой цикл.
Главная идея в том, чтобы использовать шаблон, который принимает не начало и конец последовательности, а всю последовательность целиком, и вспомогательный шаблон, который генерирует последовательность.
\begin{cppcode}
namespace EL {
...
template<class T, T begin, T end, T step = 1>
using MakeList = /*implementation*/;
}

template<From from, To to, Step step = Step{1}>
struct For : detail::ForImpl<EL::MakeList<int, from, to, step>> {};
\end{cppcode}
Несколько замечаний по коду выше:
\begin{enumerate}
\item Шаблон \verb"ForImpl" принимает всю последовательность целых индексов целиком в виде списка элементов одинакового типа и для каждого индекса из списка вызывает шаблон функции с нужным шаблонным индексом.

\item Шаблон \verb"MakeList" генерирует нужную последовательность.

\item Мы используем вместо \verb"using" директивы наследование и объявляем новую структуру для того, чтобы в ошибках компилятора у нас возникало имя структуры для которой произошла ошибка.
Мы хотим, чтобы это была \verb"For", а не просто \verb"ForImpl".
\end{enumerate}
Шаблон для генерации последовательности описан в разделе~\ref{section::EListImpl}.
Обратите внимание, что шаг от \verb"from" до \verb"to" может быть не единичный, но \verb"step = 1" является дефолтным поведением и его можно не указывать.

Теперь мы можем имплементировать сам цикл \verb"For" так
\begin{cppcode}
namespace detail {

template<class TList>
struct ForImpl {
  static_assert(ext::EL::isList<TList>, "Internal error of For Loop!");
};

template<template<class U, U...> class TList, int... Es>
struct ForImpl<TList<int, Es...>> {
  template<template<int> class F, class... Args>
  static void Do(Args&&... args) {
    (F<Es>()(std::forward<Args>(args)...), ...);
  }
  template<auto f, class... Args>
  static void Do(Args&&... args) {
    (f(Es, std::forward<Args>(args)...), ...);
  }
  template<class F, class... Args>
  static void Do(Args&&... args) {
    (F()(Es, std::forward<Args>(args)...), ...);
};
} // namespace detail
\end{cppcode}
В случае если нам передали список элементов типа \verb"int", мы вызываем нужную функцию для каждого индекса из списка с одними и теми же run-time аргументами.
И вот тут обратите внимание, как передается функция.
Как я выше писал, шаблон функции к сожалению нельзя передать в качестве аргумента, вы можете передавать лишь шаблон класса.
Так же можно передавать функцию, которая будет принимать шаблонный аргумент в качестве первого параметра или класс, у которого перегружен оператор \verb"operator()" и шаблонный аргумент так же идет в качестве первого аргумента.

Теперь обсудим, как заворачивать функцию внутрь класса.
Есть несколько способов:
\begin{enumerate}
\item Сделать статическую функцию внутри класса.
Однако, в этом случае требуется договориться об использовании единого имени, чтобы все библиотеки использовали только его.

\item Сделать \verb"operator()".
Это не требует никакого дополнительного имени.
Однако, для вызова этого метода нужно создавать временный объект класса и молиться, что компилятор это оптимизирует или это никак не скажется на производительности.
Начиная с C++23 появляется \verb"static" версия оператора \verb"operator()", но звать его ничуть не приятнее, чем создавать временный объект.
\end{enumerate}
Выше представлено три подхода:
\begin{enumerate}
\item Передаем шаблон с оператором \verb"operator()".
В этом случае индекс для цикла передается через шаблонный аргумент, а run-time аргументы через аргументы оператора \verb"operator()":
\begin{cppcode}[numbers = none]
template<int index>
struct F {
  constexpr void operator()(Arg1 arg1, Arg2, arg2) const noexcept {
    /*use index and args*/
  }
};
\end{cppcode}
Здесь \verb"Arg1" и \verb"Arg2" -- типы run-time аргументов.

\item Передаем функцию или объект класса с оператором \verb"operator()".
В этом случае все аргументы передаются как run-time аргументы, включая индекс.
Индекс идет первым аргументом:
\begin{cppcode}[numbers = none]
void f(int index, Arg1 arg1, Arg2, arg2) noexcept {
    /*use index and args*/
};
\end{cppcode}

\item Передаем класс объекта с оператором \verb"operator()".
В этом случае объект надо создать налету.
Индекс передается первым аргументом в оператор \verb"operator".
\begin{cppcode}[numbers = none]
struct F {
  constexpr void operator()(int index, Arg1 arg1, Arg2, arg2) const noexcept {
    /*use index and args*/
  }
};
\end{cppcode}
\end{enumerate}
Меня несколько смущает такое положение дел, ведь никто не мешал в STL договориться для типов использовать имя \verb"type", для статических переменных имя \verb"value".
Почему нельзя было выбрать еще одно третье имя для методов для меня загадка.
Теперь покажем, как будет выглядеть вызов цикла во всех трех случаях.
Вот как будет выглядеть код пользователя.
\begin{enumerate}
\item Будем считать, что у нас есть изначальная функция:
\begin{cppcode}
template<int index>
void function(double x, char y) {...}
\end{cppcode}
В начале создадим шаблон обертку:
\begin{cppcode}
template<int index>
struct Function {
  void operator()(double x, char y) const {
    function<index>(x, y);
  }
};
\end{cppcode}
А теперь можно звать цикл так:
\begin{cppcode}
For<From{0}, To{5}>::Do<Function>(1.2, 'c');
\end{cppcode}
Это будет равносильно вызову:
\begin{cppcode}
function<0>(1.2, 'c');
function<1>(1.2, 'c');
function<2>(1.2, 'c');
function<3>(1.2, 'c');
function<4>(1.2, 'c');
\end{cppcode}

\item Пусть теперь у нас есть функция:
\begin{cppcode}
void function(int index, double x, char y) {...}
\end{cppcode}
Тогда можно напрямую вызвать
\begin{cppcode}
For<From{0}, To{5}>::Do<function>(1.2, 'c');
\end{cppcode}

\item В этом случае так же будем считать, что у нас есть функция:
\begin{cppcode}
void function(int index, double x, char y) {...}
\end{cppcode}
Теперь ее можно завернуть в класс:
\begin{cppcode}
struct F {
  void operator()(int index, double x, char y) const {
    function(index, x, y);
};
\end{cppcode}
После чего можно вызвать кода так:
\begin{cppcode}
For<From{0}, To{5}>::Do<F>(1.2, 'c');
\end{cppcode}
\end{enumerate}
Я специально сделал не интрузивное решение, в котором я не избавляюсь от оригинальной функции, а использую ее внутри оператора \verb"operator()".
При желании этот шаблон можно усложнять и дополнять чем угодно.

\subsubsection{Вариант \texttt{ForEach}}

\paragraph{Чего мы хотим}

Мы уже умеем хранить элементы в compilte-time списках.
Это могут быть элементы одного типа или любого типа.
\begin{cppcode}
namespace EL {
template<class T, T... Es>
struct List {};
}

namespace AL {
template<auto... Es>
struct List {};
}
\end{cppcode}
Теперь хочется пробежаться по списку и применить к нему какую-то операцию.
Цикл \verb"For", который мы построили выше не дает явного доступа к имени переменной счетчика снаружи и им не очень удобно пользоваться для таких списков.
Вот пример чего бы мы хотели для случая списка из любых элементов.
\begin{cppcode}
void f() { std::cout << "f()"; }
void g() { std::cout << "g()"; }

using AList = AL::List<f, g>;

template<auto f>
struct Caller {
  void operator()(const char* Name) const {
    std::cout << Name << " ";
    f();
    std::cout <<'\n';
  }
};
int main() {
  ForEach<AList>::Do<Caller>("func");
}
\end{cppcode}
В данном случае список заполнен функциями, где каждая функция печатает свое имя.
Мы хотим вызвать по-порядку все функции и напечатать дополнительно информацию.
Вывод в этом вызове будет
\begin{cppcode}
func = f()
func = g()
\end{cppcode}
В случае списка из элементов одного типа мы бы ожидали что-то такое
\begin{cppcode}
using EList = EL::List<int, 1, 2, 3, 4>;

template<class T, T E>
struct Printer {
  void operator()(const char* Name) const {
    std::cout << Name << " = " << E << '\n';
  }
};

template<int Elem>
using PrinterInt = Printer<int, Elem>;

int main() {
  ForEach<EList>::Do<PrinterInt>("value");
}
\end{cppcode}
В этом случае вывод будет:
\begin{cppcode}
value = 1
value = 2
value = 3
value = 4
\end{cppcode}

\paragraph{Проблема с лишним аргументом}

Обратите внимание на один неприятный момент.
У нас для задания принтера надо сначала задать его тип, а потом уже элемент который мы будем печатать.
\begin{cppcode}
template<class T, T E>
struct Printer {
  void operator()(const char* Name) const {
    std::cout << Name << " = " << E << '\n';
  }
};
\end{cppcode}
Однако, метод \verb"Do" принимает шаблон вида
\begin{cppcode}
template<T E>
struct F {};
\end{cppcode}
Мы решаем эту проблему биндингом типа в псевдониме
\begin{cppcode}
template<int Elem>
using PrinterInt = Printer<int, Elem>;
\end{cppcode}
В целом можно было бы сделать еще проще и создать шаблон
\begin{cppcode}
template<auto E>
struct Printer {
  void operator()(const char* Name) const {
    std::cout << Name << " = " << E << '\n';
  }
};
\end{cppcode}
Так же можно сделать специальный шаблон-контекст, который уже заранее в себе хранит все что надо:
\begin{cppcode}
template<class T>
struct Context {
  template<T Elem>
  struct Printer {
    void operator()(const char* Name) const {
      std::cout << Name << " = " << Elem << '\n';
    }
  };
};
\end{cppcode}
Тогда вызов в цикле бы выглядел так
\begin{cppcode}
ForEach<AList>::Do<Context<int>::Printer>("value");
\end{cppcode}
У такого подхода проблема в том, что название класса \verb"Context" надо будет каждый раз придумывать разное или прятать во вспомогательное пространство имен, что раздует название <<метода>> для выполнения.

Есть еще один вариант биндинга.
Метод биндинга с помощью псевдонимов не позволяет делать биндинг налету.
Приходится вводить явно само название псевадонима.
А мы бы хотели что-то такое:
\begin{cppcode}
ForEach<AList>::Do<Print<int, >>("value");
\end{cppcode}
Но так к сожалению это не работает в плюсах.
Потому надо это делать дополнительными танцами с бубном.
Для этого заведем дополнительный класс биндер.
\begin{cppcode}
template<class T, template<class U, U> class F>
struct Bind {
  template<T Elem>
  using type = F<T, Elem>;
};
\end{cppcode}
Теперь можно сделать вызов так
\begin{cppcode}
ForEach<EList>::Do<Bind<int, Printer>::type>("value");
\end{cppcode}
В этом случае недостаток в том, что надо название биндера прятать где-то, что опять же раздует размер <<имени>> для <<метода>>.
И придется использовать служебное \verb"type".

\subsubsection{Имплементация \texttt{ForEach}}

У нас сейчас есть три типа списков:
\begin{enumerate}
\item Список элементов фиксированного типа.
\begin{cppcode}[numbers = none]
namespace EL {
template<class T, T... Es>
struct List {};
}
\end{cppcode}

\item Список произвольных элементов.
\begin{cppcode}[numbers = none]
namespace AL {
template<auto... Es>
struct List {};
}
\end{cppcode}

\item Список типов.
\begin{cppcode}[numbers = none]
namespace TL {
template<class... Ts>
struct List {};
}
\end{cppcode}
Для каждого из списков нужно сделать свою версию \verb"ForEach" цикла.
Однако, пользователь будет пользоваться едиными шаблонами снаружи:
\begin{cppcode}
template<class TList>
struct ForEach : detail::ForEachImpl<TList> {};
\end{cppcode}
Вот как будут выглядеть имплементация.
В начале общее определение, которое используется для сообщения об ошибке, так как в этом случае зовется именно он:
\begin{cppcode}
template<class TList>
struct ForEachImpl {
  static_assert(
      EL::isList<TList> || AL::isList<TList> || TL::isList<TList>,
      "The argument of ForEach MUST be either an element list or a type list!");

};
\end{cppcode}
Теперь имплементации для разных видов списков.
В начале для списка с элементами одного типа:
\begin{cppcode}
namespace detail {
template<template<class U, U...> class TList, class T, T... Es>
struct ForEachImpl<TList<T, Es...>> {
  template<template<T> class F, class... Args>
  static constexpr void Do(Args&&... args) {
    (F<Es>()(std::forward<Args>(args)...), ...);
  }
  template<class F, class... Args>
  static constexpr void Do(Args&&... args) {
    (F()(Es, std::forward<Args>(args)...), ...);
  }
  template<auto f, class... Args>
  static constexpr void Do(Args&&... args) {
    (f(Es, std::forward<Args>(args)...), ...);
  }
};
} // namespace detail
\end{cppcode}
Теперь для списка из произвольных элементов:
\begin{cppcode}
namespace detail {
template<template<auto...> class TList, auto... Es>
struct ForEachImpl<TList<Es...>> {
  template<template<auto> class F, class... Args>
  static constexpr void Do(Args&&... args) {
    (F<Es>()(std::forward<Args>(args)...), ...);
  }
  template<class F, class... Args>
  static constexpr void Do(Args&&... args) {
    (F()(Es, std::forward<Args>(args)...), ...);
  }
  template<auto F, class... Args>
  static constexpr void Do(Args&&... args) {
    (F(Es, std::forward<Args>(args)...), ...);
  }
};
} // namespace detail
\end{cppcode}
И наконец-то для списка типов:
\begin{cppcode}
namespace detail {
template<template<class...> class TList, class... Ts>
struct ForEachImpl<TList<Ts...>> {
  template<template<class> class F, class... Args>
  static constexpr void Do(Args&&... args) {
    (F<Ts>()(std::forward<Args>(args)...), ...);
  }
};
} // namespace detail
\end{cppcode}
Тут обратите внимание, что есть всего лишь один способ действовать на типах, а потому у нас всего один метод \verb"Do".
\end{enumerate}


